% page 165 du fac-simile (image p-164 du scan)
% bloc 165.1 x 242.0 mm ; marge gauche 36.9 mm
\pgc{15.240mm}{0.000mm}{
\l{\marge[78.398mm]{32.825mm}{\ornamento{11.044mm}{ornements/letroj/litero-F.png}}}
\l{\cel{53}155}
\l{}
\l{}
\l{}
\l{}
\l{}
\l{}
\l{}
\l{}
\l{}
\l{fa. (muziko).La noto quaresma di la \textquotedbl{}do\textquotedbl{}-gamo.}
\l{}
\l{\sou{fabo}. (\sou{bot.}) Planto ek la familio \textquotedbl{}leguminosi\textquotedbl{}, kun shelo grosa,}
\l{\cel{3}qua kontenas grano oblonga, uzata kom nutrivo. - L.\cel{1}\sou{vicia}}
\l{\cel{3}\sou{faba}. - FIS.}
\l{}
\l{\sou{fablo}.\cel{2}Naraco kurta, proza o versa, di qua la protagonisti esas,}
\l{\cel{3}maxim-multa-kaze, od enti imaginita, od animali, e qua uzesas}
\l{\cel{3}kom pruviva pri leciono di etiko praktikala. - DEFIS.}
\l{}
\l{\sou{fabrika}r (\sou{trans.)}\cel{2}Facar, per materii prima quin onu laboras, ula}
\l{\cel{3}kategorii de objekti destinita a livresor engrose a la komer-}
\l{\cel{3}cisti,segun procedi mashinala o mekanikala. - DEFIRS.}
\l{}
\l{\sou{facar}. (\sou{trans.)}\cel{3}Tirar kozo ek altra, adaptante lu a destineso}
\l{\cel{3}nova. - FIS.}
\l{}
\l{\sou{faceto}. Latero di objekto plura-angula. - DEFIRS.}
\l{}
\l{\sou{facio}. La parto \textquotedbl{}vizaja\textquotedbl{} di irga korpo. - EFIS.}
\l{}
\l{\sou{\textquotedbl{}facies\textquotedbl{}}.\cel{2}(\sou{cienco}). Aspekto, fizionomio.}
\l{}
\l{facila. Facebla od agebla sen peno. - EFIS.}
\l{}
\l{\sou{facinar}.\cel{2}(\sou{trans}.)I. Seduktar nerezisteble per la povo di la}
\l{\cel{3}regardo. - II. (\sou{metaf}.) Dazlar per prestijo. - DEFIS.}
\l{}
\l{facionar. (\sou{netrans.}) (\sou{pri soldato})\cel{2}Komisesir sorgor por la}
\l{\cel{3}sekureso di posteno, di institucuro publika avan qua lu stacas}
\l{\cel{3}armoze dum plura hori til remplaseso da altru. - F.}
\l{}
\l{fago.\cel{2}(\sou{bot.)}\cel{2}Arboro granda, ek la familio \textquotedbl{}amentacei\textquotedbl{}, di qua}
\l{\cel{3}la kortico esas grizatra, la foliaro densa. - L.\cel{1}\sou{fagus}. ISL.}
\l{}
\l{fagocito. (\sou{biol}.) Celulo qua povas inglobar e digestar la parti-}
\l{\cel{3}kuli neorganika ed organika.}
\l{}
\l{fagocitoso. (\sou{biol.}) Rolo di la fagocito, deskovrita da Mechnikof.}
\l{}
\l{fairo.\cel{2}I. Emanajo de kaloro e de lumo, quan produktas la kom-}
\l{\cel{3}busteso di ula korpi, quan la Antiqui konsideris kom un ek la}
\l{\cel{3}quar elementi. - II. (\sou{metaf.})\cel{2}Ardoro, saporo varma ed ecitiva}
\l{\cel{3}di la drinkaji alkoholoza ; ardoro di la psiko, di la pasioni.}
\l{\cel{3}- DE.}
}
